\section{Domain strategy}
\label{sec:domains}

Everything in this section is future work \evfut{}. No domain below has been reconstructed
under the rules stated here. The domain set is a proposal for the research program in
\cref{sec:roadmap,sec:workpackages}, not a result. It is marked once, here, rather than on
every line.

\subsection{Selection principles}

A domain is chosen for the \emph{generalization question} it answers and the \emph{human
criterion} that can test it, never for how convenient its data are. Six axes drive the choice:

\begin{itemize}[noitemsep]
  \item \textbf{Dominant knowledge type.} The methodology lenses (\cref{sec:foundations}) were
    built for different kinds of knowledge, so each type is a separate test of H2
    (\cref{sec:question}).
  \item \textbf{Documentation density.} A text-rich domain flatters reconstruction; a generality
    claim needs a text-poor domain too.
  \item \textbf{Human criterion.} Without retrospective gold, a realized-loss outcome, or fresh
    blind elicitation, a domain can debug the pipeline but never test it.
  \item \textbf{Learner data.} Any claim about learner difficulty needs it, because experts cannot
    say what is difficult for a novice \evlit{} \parencite{nathan2003expert,maries2016teaching}.
  \item \textbf{Contamination risk.} Old gold diffuses into pretraining corpora and textbooks
    \evlit{} \parencite{golchin2023time,sainz2023nlp}, so post-cutoff or fresh criteria are the
    program's clean tests.
  \item \textbf{Purpose.} A domain serves \emph{scientific} validation, \emph{commercial}
    validation, both, or neither.
\end{itemize}

Easy data alone is never a reason to include a domain.

\subsection{The domain set}

\Cref{tab:domains} lists six domains. Three carry H2's gates: D1 at the screening gate MS3, and
D3 and D4 at the H2 test MS5. D2 carries H-OSS, a separate hypothesis on the organizational
track. D5 and D6 test the learner link and a best-case ceiling; they gate nothing for H2.

\begin{table}[tbp]
  \centering
  \small
  \caption{The six domains, the generalization question each tests, its gold or human
  criterion, its role in the validation program, and its main risk. D1, D3 and D4 carry H2's
  gates; D2 carries H-OSS only; D5--D6 do not gate H2.}
  \label{tab:domains}
  \begin{tabularx}{\linewidth}{@{}p{1.7cm}XXp{2.0cm}X@{}}
    \toprule
    Domain & Generalization question & Gold / human criterion & Role & Main risk \\
    \midrule
    D1 Clinical procedural skill &
      Perceptual-motor, decision-heavy procedure: does the gap map find omitted steps and cues? &
      Retrospective CTA golds, pooled: Sullivan et al.\ cricothyrotomy task list (46 items) first, others found by the Phase~1 metadata check; NICU cue list as a candidate &
      Scientific (V2 pooled; MS3 screening) &
      Contamination of old golds; fewer than three eligible golds; dated corpora too small \\
    D2 Open-source software maintenance &
      Does artifact coverage of design rationale predict realized knowledge loss beyond authorship concentration? &
      Realized-loss outcome: blind-coded rationale-seeking issues and defect-fix commits after a truck-factor departure &
      Scientific (V3; tests H-OSS, which cannot stop or license H2) and commercial precursor
      (V7) &
      Repositories in pretraining; authorship metrics may break under agentic coding; effective $n$ = number of departures \\
    D3 Industrial fault diagnosis &
      Does the method hold where diagnostic craft dominates and text is poor? &
      Retrospective troubleshooting data (Chao \& Salvendy); PARI maintenance task list (availability unverified); prospective blind elicitation (V5) &
      Scientific (text-poor falsification target; V5, MS5) and commercial (second-wave partner) &
      PLC ledger is in-sample; low reconstruction support; technician recruitment \\
    D4 Regulatory judgment &
      Does the method work on normative-interpretive judgment, where the residual sits in hedges and exception boundaries? &
      No public itemized gold; prospective blind elicitation of practitioners (V5) &
      Scientific (V5, MS5); commercial only later &
      GDPR ledger is in-sample; practitioners may legitimately disagree \\
    D5 School mathematics &
      Does an expert-text residual relate to learner difficulty, or are they different constructs? &
      Learner response data (public misconception and response datasets) &
      Scientific only (learner link; not a calibration set) &
      Non-commercial data licenses; a null result is likely and informative \\
    D6 Introductory mechanics (Track A) &
      Best-case ceiling: the only place where expert residual, learner error, and a randomized controlled trial (RCT) of learning meet in one topic &
      Track A's own registered K0/K1/K2; sealed predictions scored after K1 is locked &
      Scientific, secondary and non-gating (V8b) &
      Contamination (Force Concept Inventory in pretraining); eligibility rule; must not change any registered rule \\
    \bottomrule
  \end{tabularx}
\end{table}

\textbf{D1.} Sullivan et al.'s cricothyrotomy task list gives 46 steps (14 clinical-knowledge, 27
action, 5 decision). Surgeons described an average of 44\% of them unprompted and 66\% when
prompted \evlit{} \parencite{sullivan2014use}. A NICU cue list also exists, but whether it can be
reused is unverified \evlit{} \parencite{crandall1993critical}. One gold cannot screen H2, so V2
pools several published CTA golds; the Phase~1 metadata check decides how many exist. D1 is the
best-documented perceptual-motor procedure with a published, itemized criterion.

\textbf{D2.} Developer departure gives a realized-loss outcome that needs no participant or consent
to measure retrospectively \evlit{} \parencite{avelino2016novel,avelino2019abandonment,rigby2016quantifying,nassif2017revisiting}.
It tests H-OSS: the gap map, adapted to code and commit evidence, predicts where knowledge is lost
when developers leave. The adaptation needs a Git adapter and code lenses (for example WHY on
commits with no linked rationale). Both are built on the organizational track, and the code lenses
are frozen only after the adapter exists (\cref{sec:workpackages}). D2 is the commercial
precursor, but it tests a different construct from H2 and can neither stop nor license it.

\textbf{D3.} Chao and Salvendy's troubleshooting study \evlit{} \parencite{chao1994percentage} and
the PARI task list \evlit{} \parencite{hall1995procedural} anchor a domain where diagnostic craft,
not declarative content, carries the knowledge. Included to test generality on a text-poor domain
(\cref{sec:question}); PLC/SCADA text serves only as its development subdomain.

\textbf{D4.} No public itemized gold exists for normative-interpretive judgment, so this domain
tests whether the residual sits in hedges and exception boundaries a fresh, blind elicitation can
reveal, following work-as-imagined versus work-as-done \evlit{} \parencite{hollnagel2018safety}.
Included because D1--D3 do not touch normative judgment; GDPR/DPIA text is only its development
subdomain.

\textbf{D5.} Public misconception and response datasets give rich, consented learner data with no
participant recruited. Included only to ask whether an expert-text residual predicts learner
difficulty, separable from RQ-B; a null result here is informative and does not falsify H2.

\textbf{D6.} Track A (\cref{sec:validation}) is the one place expert residual, learner error, and a
learning RCT already meet in one design. Included as a secondary, non-gating comparison only,
through E-SEAL; it drives no decision in this program.

\subsection{Domains considered and not chosen}

Convenience of data was deliberately excluded as a selection reason. Mathematics has the easiest
public data of any domain considered. But it measures learner difficulty, not the human knowledge
residual, so it can never be a calibration set for the other five domains. Customer-support
diagnostics and legal contract review were set aside for this round: no realized-loss outcome or
itemized gold could be identified for them. They remain candidates for the new non-gold domain in
the Phase~0 closure study (V1).

\subsection{Binding rules}

Three rules shape the set and must not be read as recommendations.

\begin{enumerate}
  \item \textbf{Physics is the frozen Track A domain.} Introductory mechanics enters the program
    only through E-SEAL: sealed, script-generated predictions, non-gating. NOW runs exclude
    conservation items entirely. Anyone who has seen reconstructed conservation content is
    ineligible for any Track A human role -- interviewer, corroborating coder, K1 absent-operation
    coder, or importance rater.
  \item \textbf{Held-out gold domains are not reconstructed first.} No exploratory, debugging, or
    tuning run touches a gold domain before its configuration is hash-frozen and the study is
    pre-registered. The first reconstruction of a gold domain \emph{is} the test run, on a corpus
    dated before the gold existed.
  \item \textbf{Development domains stay development domains.} PLC and GDPR tuned the gap map's
    lexicons in-sample. They may be used only for debugging, for the Phase~0 closure study (V1),
    and for a non-analyzed V5 pilot. They never carry a gate reading.
\end{enumerate}

The clearest tension this creates is in choosing the domain for V5, the prospective test of H2.
One earlier note suggested PLC because it is already reconstructed. That conflicts with rule 3:
the lexicons were tuned on the PLC ledger, so any PLC effect would be optimistic.
The recommended resolution, in order of merit:

\begin{itemize}[noitemsep]
  \item \textbf{V5 pilot in PLC, V5 main in two held-out subdomains (Recommended).} The pilot
    is not analyzed for effect, so in-sample tuning cannot bias it. The main study runs once,
    after the freeze, on a fresh maintenance-diagnosis task (D3) and a held-out regulation (D4).
    The extra cost is one reconstruction and one area partition per domain, small next to the
    human study.
  \item V5 main in PLC and GDPR directly. This saves one partition each. But the prospective
    test of H2 would then test a configuration on its own tuning data, which weakens any positive
    result and cannot be repaired later.
\end{itemize}
What would change the recommendation: no reachable expert pool for a fresh maintenance
subdomain. Then keep PLC but re-tune nothing, state the in-sample status in the pre-registration,
and treat the PLC result as supporting evidence only, never as the reading that tests H2.
